% Auto-generated by report/make_tables.py -- do not edit by hand.

\begin{table*}[t]
\centering\small
\setlength{\tabcolsep}{5pt}
\begin{tabular}{lrrrrrcrrrr}
\toprule
& \multicolumn{5}{c}{\textbf{Quality (greedy decoding)}} & & \multicolumn{4}{c}{\textbf{Cost}} \\
\cmidrule(lr){2-6}\cmidrule(lr){8-11}
\textbf{Cfg} & \textbf{Bit \%} & \textbf{Seq \%} & \textbf{Lev.} & \textbf{BLEU} & \textbf{R-L} & & \textbf{Par (M)} & \textbf{s/ep} & \textbf{Peak MB} & \textbf{Dec (s)} \\
\midrule
C1 & 97.17 & 77.20 & 0.44 & 90.34 & 95.10 & & 7.62 & 143.0 & 372 & 6.6 \\
C2 & \textbf{97.61} & \textbf{81.20} & \textbf{0.35} & \textbf{92.27} & \textbf{96.06} & & 7.62 & 153.8 & 372 & 6.6 \\
C3 & 96.81 & 73.80 & 0.54 & 88.54 & 94.21 & & 6.73 & 147.2 & 360 & 6.1 \\
C4 & 97.33 & 78.80 & 0.42 & 90.91 & 95.42 & & 7.61 & 131.9 & 350 & 5.7 \\
C5 & 97.02 & 70.00 & 1.00 & 82.80 & 91.23 & & 6.71 & 275.7 & 523 & 27.6 \\
\bottomrule
\end{tabular}
\caption{Test-set results (25 epochs each); best per column in \textbf{bold}. Levenshtein is mean edit distance per segment, lower is better. Cost measured on one RTX 3050 Laptop GPU (4\,GB, bfloat16), each configuration in its own process; `Dec.'\ decodes 1{,}000 segments.}
\label{tab:quality}
\end{table*}

